\textbf{Phần 1. Trắc nghiệm lựa chọn}
\setcounterpageref{ex}{0}

\begin{ex} %[3D2V63]
Tính giá trị của phép tính: $45\,280 + 32\,415$.
\choice
{$77\,595$}
{$77\,795$}
{\True $77\,695$}
{$76\,695$}
\loigiai{
$45\,280 + 32\,415 = 77\,695$.
\textbf{Đáp án của câu hỏi là C}
}
\end{ex}

\begin{ex} %[3D2V63]
Một nhà máy tháng thứ nhất sản xuất được $28\,500$ sản phẩm, tháng thứ hai sản xuất được $31\,200$ sản phẩm. Cả hai tháng nhà máy sản xuất được bao nhiêu sản phẩm?
\choice
{\True $59\,700$ sản phẩm}
{$59\,500$ sản phẩm}
{$60\,000$ sản phẩm}
{$58\,700$ sản phẩm}
\loigiai{
$28\,500 + 31\,200 = 59\,700$ sản phẩm.
\textbf{Đáp án của câu hỏi là A}
}
\end{ex}

\begin{ex} %[3D2V63]
Tìm số tự nhiên $x$, biết: $x - 24\,500 = 35\,200$.
\choice
{$58\,700$}
{$60\,700$}
{$59\,000$}
{\True $59\,700$}
\loigiai{
$x = 35\,200 + 24\,500 = 59\,700$.
\textbf{Đáp án của câu hỏi là D}
}
\end{ex}

\begin{ex} %[3D2V63]
Tính nhanh giá trị biểu thức: $25\,000 + 18\,400 + 15\,000 + 1600$.
\choice
{$58\,000$}
{\True $60\,000$}
{$62\,000$}
{$59\,000$}
\loigiai{
$(25\,000 + 15\,000) + (18\,400 + 1600) = 40\,000 + 20\,000 = 60\,000$.
\textbf{Đáp án của câu hỏi là B}
}
\end{ex}

\begin{ex} %[3D2V63]
Một huyện năm ngoái trồng được $34\,500$ cây xanh, năm nay trồng được nhiều hơn năm ngoái $5200$ cây xanh. Cả hai năm huyện đó trồng được bao nhiêu cây xanh?
\choice
{$39\,700$ cây}
{$73\,200$ cây}
{\True $74\,200$ cây}
{$75\,000$ cây}
\loigiai{
Năm nay trồng: $34\,500 + 5200 = 39\,700$ cây. Cả hai năm: $34\,500 + 39\,700 = 74\,200$ cây.
\textbf{Đáp án của câu hỏi là C}
}
\end{ex}

\begin{ex} %[3D2V63]
Tìm tổng của số lớn nhất có năm chữ số và số bé nhất có năm chữ số.
\choice
{\True $109\,999$}
{$110\,000$}
{$99\,999$}
{$100\,000$}
\loigiai{
Số lớn nhất: 99 999. Số bé nhất: 10 000. Tổng: $99\,999 + 10\,000 = 109\,999$.
\textbf{Đáp án của câu hỏi là A}
}
\end{ex}

\begin{ex} %[3D2V63]
Một kho hàng có $42\,000\text{ kg}$ gạo. Người ta nhập thêm đợt 1 là $18\,500\text{ kg}$ và đợt 2 là $15\,500\text{ kg}$. Trong kho có tất cả bao nhiêu ki-lô-gam gạo?
\choice
{$75\,000\text{ kg}$}
{$77\,000\text{ kg}$}
{$74\,000\text{ kg}$}
{\True $76\,000\text{ kg}$}
\loigiai{
Tổng gạo: $42\,000 + 18\,500 + 15\,500 = 76\,000\text{ (kg)}$.
\textbf{Đáp án của câu hỏi là D}
}
\end{ex}

\begin{ex} %[3D2V63]
Điền số thích hợp: $56\,340 + \dots = 80\,000$.
\choice
{$24\,660$}
{\True $23\,660$}
{$23\,760$}
{$22\,660$}
\loigiai{
$80\,000 - 56\,340 = 23\,660$.
\textbf{Đáp án của câu hỏi là B}
}
\end{ex}

\begin{ex} %[3D2V63]
Một gia đình mua một chiếc tivi giá $12\,500\,000$ đồng và một máy giặt giá $8\,200\,000$ đồng. Tổng số tiền gia đình đó phải trả là
\choice
{$20\,500\,000$ đồng}
{$21\,000\,000$ đồng}
{\True $20\,700\,000$ đồng}
{$19\,700\,000$ đồng}
\loigiai{
$12\,500\,000 + 8\,200\,000 = 20\,700\,000$ đồng.
\textbf{Đáp án của câu hỏi là C}
}
\end{ex}

\begin{ex} %[3D2V63]
Nếu số hạng thứ nhất tăng thêm 1500 đơn vị và số hạng thứ hai tăng thêm 2500 đơn vị thì tổng sẽ tăng thêm bao nhiêu?
\choice
{\True $4000$ đơn vị}
{$3000$ đơn vị}
{$5000$ đơn vị}
{$3500$ đơn vị}
\loigiai{
Tổng tăng: $1500 + 2500 = 4000$ đơn vị.
\textbf{Đáp án của câu hỏi là A}
}
\end{ex}

\begin{ex} %[3D2V63]
Tính giá trị: $48\,500 + 12\,500 + 20\,000$.
\choice
{$80\,000$}
{$82\,000$}
{$79\,000$}
{\True $81\,000$}
\loigiai{
$48\,500 + 12\,500 = 61\,000$; $61\,000 + 20\,000 = 81\,000$.
\textbf{Đáp án của câu hỏi là D}
}
\end{ex}

\begin{ex} %[3D2V63]
Một công ty thu nhập quý Một là $45\,000\,000$ đồng, quý Hai thu nhập nhiều hơn quý Một là $8\,000\,000$ đồng. Thu nhập cả hai quý của công ty là
\choice
{$53\,000\,000$ đồng}
{\True $98\,000\,000$ đồng}
{$90\,000\,000$ đồng}
{$100\,000\,000$ đồng}
\loigiai{
Quý Hai: $45 + 8 = 53$ triệu. Cả hai quý: $45 + 53 = 98$ triệu đồng.
\textbf{Đáp án của câu hỏi là B}
}
\end{ex}

\textbf{Phần 2. Trắc nghiệm đúng sai}
\setcounterpageref{ex}{0}

\begin{ex} %[3D2V63]
Xét tính đúng sai của phép cộng các số trong phạm vi 100 000:
\choiceTF
{\True Cộng lần lượt từ phải sang trái (từ hàng đơn vị).}
{$45\,000 + 55\,000 = 90\,000$.}
{\True Khi cộng số có 5 chữ số, nếu có nhớ thì ta thêm vào hàng liền kề bên trái.}
{$62\,300 + 14\,700 = 76\,000$.}
\loigiai{
\begin{itemize}
    \item Ý a) Đúng: Quy tắc chuẩn.
    \item Ý b) Sai: Sai vì $45\,000 + 55\,000 = 100\,000$.
    \item Ý c) Đúng: Quy tắc cộng có nhớ.
    \item Ý d) Sai: Sai vì $62\,300 + 14\,700 = 77\,000$.
\end{itemize}
\begin{itemize}
    \item a) \textbf{Đúng}.
    \item b) \textbf{Sai}.
    \item c) \textbf{Đúng}.
    \item d) \textbf{Sai}.
\end{itemize}
}
\end{ex}

\begin{ex} %[3D2V63]
Một nông trường có hai trại nuôi bò: trại 1 có 18 500 con, trại 2 có 21 500 con:
\choiceTF
{\True Cả hai trại có tất cả 40 000 con bò.}
{\True Trại 2 nhiều hơn trại 1 là 3000 con bò.}
{Tổng số bò làm tròn đến hàng chục nghìn là 50 000 con.}
{\True Nếu nhập thêm 5000 con bò nữa thì toàn nông trường có 45 000 con bò.}
\loigiai{
\begin{itemize}
    \item Ý a) Đúng: $18\,500 + 21\,500 = 40\,000$ con.
    \item Ý b) Đúng: $21\,500 - 18\,500 = 3000$ con.
    \item Ý c) Sai: Sai vì 40 000 là số tròn chục nghìn rồi.
    \item Ý d) Đúng: $40\,000 + 5000 = 45\,000$ con.
\end{itemize}
\begin{itemize}
    \item a) \textbf{Đúng}.
    \item b) \textbf{Đúng}.
    \item c) \textbf{Sai}.
    \item d) \textbf{Đúng}.
\end{itemize}
}
\end{ex}

\textbf{Phần 3. Trả lời ngắn}
\setcounterpageref{ex}{0}

\begin{ex} %[3D2V63]
Tính kết quả phép cộng: $54\,628 + 27\,345$.
\shortans{81973}
\loigiai{
\textbf{Đáp số:} $81973$\\
$54\,628 + 27\,345 = 81\,973$.
}
\end{ex}

\begin{ex} %[3D2V63]
Tìm $x$, biết: $x - 34\,200 = 45\,800$.
\shortans{80000}
\loigiai{
\textbf{Đáp số:} $80000$\\
$x = 45\,800 + 34\,200 = 80\,000$.
}
\end{ex}

\begin{ex} %[3D2V63]
Tính nhanh giá trị: $36\,000 + 28\,500 + 14\,000 + 11\,500$.
\shortans{90000}
\loigiai{
\textbf{Đáp số:} $90000$\\
$(36\,000 + 14\,000) + (28\,500 + 11\,500) = 50\,000 + 40\,000 = 90\,000$.
}
\end{ex}

\begin{ex} %[3D2V63]
Một kho hàng buổi sáng xuất $15\,400\text{ kg}$ đường, buổi chiều xuất $18\,600\text{ kg}$ đường. Hỏi cả ngày kho xuất bao nhiêu ki-lô-gam đường?
\shortans{34000}
\loigiai{
\textbf{Đáp số:} $34000$\\
$15\,400 + 18\,600 = 34\,000\text{ (kg)}$.
}
\end{ex}

\begin{ex} %[3D2V63]
Tính tổng của số tròn chục nghìn lớn nhất có 5 chữ số và số tròn nghìn bé nhất có 5 chữ số.
\shortans{100000}
\loigiai{
\textbf{Đáp số:} $100000$\\
$90\,000 + 10\,000 = 100\,000$.
}
\end{ex}

\begin{ex} %[3D2V63]
Một nhà máy dệt được $42\,500\text{ m}$ vải hoa và $38\,500\text{ m}$ vải trắng. Hỏi nhà máy dệt được tất cả bao nhiêu mét vải?
\shortans{81000}
\loigiai{
\textbf{Đáp số:} $81000$\\
$42\,500 + 38\,500 = 81\,000\text{ (m)}$.
}
\end{ex}
